% ---------------------------------------------------------------
%  Section II -- Related Work
% ---------------------------------------------------------------

\subsection{Federated Learning and Privacy-Aware Mobility}

Federated learning~\cite{mcmahan2017fedavg} enables model training across
distributed data holders without transmitting raw records to a central server.
McMahan et al.\ introduced FedAvg, in which clients execute $\tau$ local
stochastic gradient descent steps and the server aggregates parameter
updates weighted by local sample size; this mechanism is the conceptual
basis for the IRIDIUM FL pathway.
Li et al.\ \cite{li2020fedprox} showed that heterogeneous client drift can
be mitigated by adding a proximal regularization term
$\mu/2\,\norm{\theta - \theta^{(r)}}^2$ to each client's objective.
Li et al.\ \cite{li2020fedavg_convergence} established convergence bounds
for FedAvg under non-IID data distributions, which is a practically
critical condition when mobility data are partitioned by operator or
district: for non-convex objectives, FedAvg converges at rate
$\mathcal{O}(1/\sqrt{TK})$ under bounded gradient variance.
Kairouz et al.\ \cite{kairouz2021advances} provide a comprehensive survey
covering Byzantine robustness, communication compression, and differential
privacy, all relevant to a multi-operator Azerbaijani deployment.
Lim et al.\ \cite{lim2020federated_edge} survey FL in mobile edge networks,
highlighting that latency constraints and device heterogeneity impose strict
bounds on communication rounds.
IRIDIUM positions its FL component as a fully specified architectural
extension point awaiting local training node deployment.

\subsection{Spatio-Temporal Traffic Forecasting}

Short-horizon traffic forecasting has been substantially advanced by
graph-structured deep learning.
Li et al.\ \cite{li2018dcrnn} proposed the diffusion convolutional recurrent
network (DCRNN), which models traffic as a bidirectional diffusion process
on the road graph and outperforms recurrent baselines on METR-LA and
PeMS-BAY by leveraging spatial correlations.
Yu et al.\ \cite{yu2018stgcn} introduced spatio-temporal graph convolutional
networks (ST-GCN) using purely convolutional temporal operations for
faster training.
Graph WaveNet~\cite{wu2019graphwavenet} extended the framework with
adaptive, self-learned adjacency matrices that capture non-local dependencies
absent from static proximity graphs.
Guo et al.\ \cite{guo2019astgcn} added attention mechanisms over both spatial
and temporal dimensions, yielding attention-based ST-GCN (ASTGCN) with
improved interpretability on PEMS datasets.
Jiang and Luo~\cite{jiang2022gnn_survey} survey graph neural networks for
traffic forecasting comprehensively.
The ST-GNN extension point in IRIDIUM follows Eq.~(\ref{eq:forecast})
and Eq.~(\ref{eq:gcn}); a heuristic baseline is active until a historical
observation store of sufficient length is assembled.

\subsection{Urban Digital Twins and Transport Platforms}

The concept of a city-scale digital twin is gaining traction in research
and practice.
Ferre-Bigorra et al.\ \cite{ferrebigorra2022urban_digital_twins} analyze
adoption barriers including data governance and interoperability, which
directly motivate the provenance-and-status design of IRIDIUM.
Wan et al.\ \cite{wan2024urban_modelling_to_city_digital_twins} trace the
transition from offline urban modelling to continuously updated digital
twins with live sensor fusion.
Weil et al.\ \cite{weil2023urban_digital_twin_challenges} systematically
review challenges and sustainability requirements.
Nag et al.\ \cite{nag2025digital_twins_transport_planning} and
Yan et al.\ \cite{yan2023digital_twin_transport_infra} survey digital
twin applications in transport planning and infrastructure management.
A key distinction of IRIDIUM is that the twin is assembled dynamically from
real configured sources at every request, rather than frozen at a
batch-update boundary.

\subsection{Multimodal Routing and Accessibility}

Delling et al.\ \cite{delling2015raptor} introduced the RAPTOR algorithm
for round-based public-transit routing, achieving practical query times
on city-scale networks through range-based Dijkstra extensions.
Pfertner et al.\ \cite{pfertner2023multimodal_accessibility} describe an
open-source methodology for multimodal accessibility analysis using
OpenTripPlanner and PostGIS.
Graff et al.\ \cite{graff2024nomad} construct a routable, multi-cost,
time-dependent network suitable for combined road and transit optimization.
Gil~\cite{gil2015multimodal_osm_accessibility} builds a multimodal network
directly from OpenStreetMap for sustainable urban accessibility analysis.
IRIDIUM generalizes the routing objective to the weighted sum in
Eq.~(\ref{eq:routing}), covering travel time, monetary cost, carbon
emissions, and transfer penalties; full OpenTripPlanner integration is
deferred until GTFS feeds are available.

\subsection{Mobility Equity Measurement}

Karner et al.\ \cite{karner2025transportation_equity} discuss advances and
methodological pitfalls in measuring transportation equity, noting that
composite indices must account for indicator selection bias and distributional
assumptions.
Guo et al.\ \cite{guo2020transportation_equity_overview} provide a systematic
overview spanning accessibility, emissions, and safety dimensions.
Martens et al.\ \cite{martens2019measuring_transport_equity_metrics} and
Lucas et al.\ \cite{lucas2019measuring_transport_equity_book} define
key components and metrics for transport equity measurement.
The IRIDIUM Mobility Equity Score (Eq.~\ref{eq:mes}) is a district-level
composite of normalized indicators reported as a derived analytic index and
not as an official government statistic.

\subsection{Anomaly Detection in Transportation}

Pramanik et al.\ \cite{pramanik2022traffic_anomaly_video} employ
spatio-temporal rough-fuzzy granulation for video-based anomaly detection
in road surveillance.
Hassan et al.\ \cite{hassan2019spatiotemporal_anomaly_its} address
spatio-temporal detection in intelligent transportation systems using
statistical and machine-learning methods.
IRIDIUM adopts a rule-based composite severity score
(Eq.~\ref{eq:anomaly_severity}) that fuses rolling z-scores over speed
deviations with event-proximity and adverse-weather signals.

\subsection{Comparison with Related Systems}

Table~\ref{tab:related_systems} positions IRIDIUM against representative
urban mobility and traffic forecasting platforms across six dimensions.
The key differentiators of IRIDIUM are: (i) the strict no-synthetic-substitution
policy enforced through a machine-readable status protocol; (ii) the
FL architectural pathway for multi-operator privacy preservation; and
(iii) the explicit targeting of a data-scarce, data-governance-constrained
environment in the South Caucasus.

\begin{table*}[t]
\caption{Comparison of IRIDIUM with representative urban mobility and traffic
         forecasting platforms. Checkmarks indicate supported or planned
         capabilities; dashes indicate that the property is absent or
         not applicable. \textit{Arch.} denotes an architected capability
         not yet evaluated experimentally.}
\label{tab:related_systems}
\centering
\renewcommand{\arraystretch}{1.20}
\begin{tabular}{@{}lcccccccc@{}}
\toprule
\textbf{Platform} &
\textbf{RT} &
\textbf{FL} &
\textbf{Prov.} &
\textbf{OSS} &
\textbf{Multi.} &
\textbf{Equity} &
\textbf{Region} &
\textbf{Open data} \\
\midrule
IRIDIUM (this work)
  & \checkmark & \checkmark$^*$ & \checkmark & \checkmark
  & \checkmark & \checkmark & AZ & Yes \\
DCRNN~\cite{li2018dcrnn}
  & \checkmark & -- & -- & \checkmark
  & -- & -- & US & Curated \\
Graph WaveNet~\cite{wu2019graphwavenet}
  & \checkmark & -- & -- & \checkmark
  & -- & -- & US & Curated \\
ASTGCN~\cite{guo2019astgcn}
  & \checkmark & -- & -- & \checkmark
  & -- & -- & US & Curated \\
Urban DT~\cite{nag2025digital_twins_transport_planning}
  & Varies & -- & Partial & Varies
  & Varies & -- & Various & Varies \\
OTP / Valhalla
  & Partial & -- & -- & \checkmark
  & \checkmark & -- & Global & GTFS/OSM \\
FedMob~\cite{lim2020federated_edge}
  & \checkmark & \checkmark & -- & --
  & -- & -- & Various & Private \\
\bottomrule
\multicolumn{9}{@{}l@{}}{\footnotesize RT: real-time; FL: federated learning ($^*$arch.); Prov.: provenance; OSS: open-source; Multi.: multimodal; AZ: Azerbaijan.}
\end{tabular}
\end{table*}

\subsection{Data Standards and Open Geospatial Infrastructure}

GTFS~\cite{gtfs_overview} and GTFS Realtime~\cite{gtfs_realtime_overview}
define the transit data standards that IRIDIUM targets for integration.
OpenStreetMap provides the base road-network topology; Klinkhardt et al.\
\cite{klinkhardt2023osm_poi_quality} assess OSM point-of-interest quality,
while Moradi et al.\ \cite{moradi2022osm_road_quality} analyze road-network
quality indicators relevant to routing accuracy.
IRIDIUM fetches the Geofabrik Azerbaijan extract, stores the SHA-256 file
checksum in a manifest table, and records full provenance for
reproducibility.
